\section{Evaluation}

\paragraph{Protocol.}
We adapt V-JEPA 2.1 at 2B and 1B parameters on a stratified 10k-clip set, and the 1B model also on the full corpus. FactorJEPA is compared with the frozen backbone, full fine-tuning, linear probing then fine-tuning \cite{kumar2022finetuning}, continual self-supervised learning (plain continued pretraining with the V-JEPA objective), LoRA \cite{hu2022lora}, DoRA \cite{liu2024dora}, and Auto-RGN \cite{lee2023surgical}, all under identical clips, masking, optimization steps, and splits. Every other trained method starts from that continued-pretraining checkpoint, which saw the training clips only, so no method sees a validation or test clip; the frozen backbone is V-JEPA 2.1 as released. Test clips are held out at the clip level (1,825 in the 10k regime, 23,106 in the full regime), so a source video can contribute clips to both training and test (Section~\ref{sec:scope}). Every margin is a paired 95\% BCa bootstrap interval over test clips.

\paragraph{Diagnostics of predictive structure.}
\emph{Future-frame L1}: error between predicted and target-encoder embeddings on the token blocks hidden by V-JEPA's mask generator. \emph{Forward L1}: the same error when the second half of a clip is predicted from its first half. \emph{Mask-ratio slope}: growth of that error as the hidden fraction rises from 0.3 to 0.9. \emph{Motion cosine}: how much closer a test clip's mean-pooled embedding lies to clips of its own motion class than to clips of other classes, with no training; the 16 motion classes come from optical flow (speed quartile by dominant direction). Nine further diagnostics probe temporal coherence and semantics (Table~\ref{tab:headline}). Small heads trained on the frozen embedding measure temporal order (pick the true frame order among permutations), pace (recognise the playback speed), action top-1 (an attentive probe for the motion class), and arrow-of-time (forward versus reversed playback). TCC $\tau$ and TCC cycle error test, with no training, whether matching moments of two clips align. With the predictor fed its own outputs, exposure bias is the error gap to feeding it the true frames, rollout drift is the error growth per step, and L1-vs-$\Delta t$ is the error growth with the prediction horizon.

\begin{table}[t]
\centering
\small
\setlength{\tabcolsep}{2pt}
\begin{tabular}{lrrrrrr}
\toprule
 & \multicolumn{2}{c}{2B$\cdot$10k} & \multicolumn{2}{c}{1B$\cdot$10k} & \multicolumn{2}{c}{1B$\cdot$115k} \\
\cmidrule(lr){2-3}\cmidrule(lr){4-5}\cmidrule(lr){6-7}
Diagnostic & Ours & Rival & Ours & Rival & Ours & Rival \\
\midrule
future-frame L1 ($\downarrow$) & \textbf{0.496} & 0.504\textsuperscript{L} & \textbf{0.590} & 0.596\textsuperscript{L} & \textbf{0.573} & 0.589\textsuperscript{L} \\
forward L1 ($\downarrow$) & \textbf{0.529} & 0.533\textsuperscript{L} & \textbf{0.614} & 0.617\textsuperscript{L} & \textbf{0.597} & 0.607\textsuperscript{L} \\
mask-ratio slope ($\downarrow$) & \textbf{0.053} & 0.054\textsuperscript{L} & \textbf{0.052} & 0.053\textsuperscript{D} & \textbf{0.047} & 0.058\textsuperscript{L} \\
motion-cos ($\uparrow$) & 0.116 & \textbf{0.244}\textsuperscript{F} & 0.048 & \textbf{0.090}\textsuperscript{F} & \textbf{0.097} & 0.078\textsuperscript{L} \\
\midrule
temporal-order ($\uparrow$) & 0.944 & \textbf{0.947}\textsuperscript{Z} & 0.925 & \textbf{0.927}\textsuperscript{C} & \textbf{0.918} & 0.915\textsuperscript{Z} \\
pace ($\uparrow$) & 0.719 & \textbf{0.724}\textsuperscript{L} & 0.711 & \textbf{0.729}\textsuperscript{Z} & 0.703 & 0.703\textsuperscript{L} \\
action top-1 ($\uparrow$) & 0.500 & \textbf{0.528}\textsuperscript{C} & 0.483 & \textbf{0.509}\textsuperscript{F} & 0.567 & \textbf{0.572}\textsuperscript{L} \\
arrow-of-time ($\uparrow$) & 0.935 & \textbf{0.965}\textsuperscript{Z} & 0.957 & \textbf{0.965}\textsuperscript{Z} & 0.955 & \textbf{0.975}\textsuperscript{Z} \\
TCC $\tau$ ($\uparrow$) & 0.274 & \textbf{0.292}\textsuperscript{Z} & 0.287 & \textbf{0.334}\textsuperscript{Z} & 0.380 & \textbf{0.400}\textsuperscript{Z} \\
exposure-bias ($\downarrow$) & 0.048 & \textbf{0.038}\textsuperscript{C} & \textbf{0.034} & 0.036\textsuperscript{D} & 0.038 & \textbf{0.035}\textsuperscript{L} \\
rollout drift ($\downarrow$) & 0.011 & \textbf{0.008}\textsuperscript{C} & 0.005 & 0.005\textsuperscript{D} & 0.007 & \textbf{0.006}\textsuperscript{L} \\
L1-vs-$\Delta t$ ($\downarrow$) & 0.008 & \textbf{0.006}\textsuperscript{C} & 0.009 & \textbf{0.007}\textsuperscript{Z} & 0.009 & \textbf{0.007}\textsuperscript{Z} \\
TCC cycle ($\downarrow$) & 2.249 & \textbf{2.111}\textsuperscript{Z} & 2.416 & \textbf{2.195}\textsuperscript{Z} & 2.806 & \textbf{2.132}\textsuperscript{Z} \\
\bottomrule
\end{tabular}
\par\smallskip
{\small\raggedright
Rival superscript: L LoRA, D DoRA, F full fine-tuning, C continual self-supervised learning (plain continued pretraining), Z frozen backbone. At 115k only LoRA and the frozen backbone were trained as rivals.\par}
\caption{One fixed FactorJEPA model (Ours, identical in every cell) against its strongest rival on 13 diagnostics. The upper four rows measure predictive structure; the lower nine measure temporal coherence and semantics, where the rival is stronger on most rows. The rival is chosen per diagnostic on the same test clips, which favours the rivals. Bold: the better of each pair (neither when equal).}
\label{tab:headline}
\end{table}

\begin{figure}[!tb]
\centering
\includegraphics[width=0.97\columnwidth]{forest_115k_col.png}
\caption{Separation of the fixed FactorJEPA model from the strongest rival at 1B on the full corpus, in widths of the paired 95\% BCa interval (symmetric-log axis; dashed line: one width). Green: FactorJEPA ahead by at least one width; orange: ahead within the interval; red: the rival ahead.}
\label{fig:forest}
\end{figure}

\paragraph{Findings.}
With limited data (Table~\ref{tab:headline}), FactorJEPA separates from the best rival on Future-frame L1 and Forward L1 at both scales and on Mask-ratio slope at 1B, while Motion cosine favors full fine-tuning. After full 115k-clip training, against the two rivals trained at that scale (LoRA and the frozen backbone), FactorJEPA separates on all four diagnostics, Motion cosine included (Figure~\ref{fig:forest}), by 13.9 (Forward L1) to 43.3 (Mask-ratio slope) interval widths, a 2.7\% lower Future-frame L1 and a 19\% flatter Mask-ratio slope than LoRA. At 115k the Motion cosine lead is over LoRA and the frozen backbone; full fine-tuning, which leads on Motion cosine at 10k, was not trained at 115k. Method rankings on the four diagnostics replicate from the 2B to the 1B backbone (Spearman $\rho \geq 0.895$). The full-data evaluation also shows a trade-off in temporal coherence: below the rule, the best rival leads at 115k on six diagnostics (rollout drift, exposure bias, L1-versus-$\Delta t$ decay, arrow-of-time, and the two TCC measures) by 8.3 to 146.5 interval widths. FactorJEPA's strength is one-step prediction of crowded scenes, while the rivals keep the advantage on the coherence of long predicted futures.

\paragraph{Why the rivals lead on coherence.} The coherence diagnostics reward temporal cues learned from general video. Of the 23 cells in the lower nine rows of Table~\ref{tab:headline} that a rival wins, 13 go to the frozen backbone, which leads on arrow-of-time and both TCC measures in all three regimes, and 5 to plain continued pretraining, the starting point of every other trained method. Adapting to the corpus at all costs some of these cues: continued pretraining already lowers arrow-of-time accuracy at 2B from 0.965 to 0.952. On the 10k test clips at 2B, the staged schedule on raw clips costs more (arrow-of-time 0.935; exposure bias from 0.038 to 0.045), and the curriculum's clips recover part of the TCC loss (Kendall $\tau$ from 0.252 to 0.274) while raising exposure bias to 0.048; the 95\% intervals before and after each of these steps do not overlap. At 1B the staged schedule shows no such cost (exposure bias 0.037, 0.035, and 0.034 for plain, RAW, and Ours).

\begin{table}[t]
\centering
\small
\setlength{\tabcolsep}{3pt}
\begin{tabular}{llrrrr}
\toprule
 & Adaptation & Fut.\ L1$\downarrow$ & Fwd.\ L1$\downarrow$ & Mask$\downarrow$ & Motion$\uparrow$ \\
\midrule
2B & plain & 0.533 & 0.558 & 0.0625 & \textbf{0.151} \\
 & RAW & 0.497 & \textbf{0.528} & 0.0565 & 0.133 \\
 & Ours & \textbf{0.496} & 0.529 & \textbf{0.0531} & 0.116 \\
\midrule
1B & plain & 0.607 & 0.628 & 0.0586 & 0.054 \\
 & RAW & \textbf{0.587} & \textbf{0.611} & 0.0528 & \textbf{0.060} \\
 & Ours & 0.590 & 0.614 & \textbf{0.0515} & 0.048 \\
\bottomrule
\end{tabular}
\caption{Where the gain comes from, 10k clips, 1,825 test clips. Plain: continued pretraining, the starting point of both other rows. RAW: the same staged adaptation as Ours on raw clips only. Ours: the staged adaptation on the curriculum's layout, agent, and interaction clips. Every RAW-minus-plain and Ours-minus-RAW difference excludes zero at 95\% under both clip and source-video resampling (127 videos), except the Ours-minus-RAW difference in Future-frame L1 at 2B.}
\label{tab:ladder}
\end{table}

\paragraph{Where the gain comes from.} Table~\ref{tab:ladder} separates two contributions. The staged adaptation is the main source of the predictive gain: on raw clips alone it lowers Future-frame L1 by 0.036 (2B) and 0.020 (1B) below plain continued pretraining, and it already beats LoRA on Future-frame and Forward L1 at both scales. The automatically built layout, agent, and interaction clips add a specific property, robustness to missing evidence: they flatten the Mask-ratio slope by \nLadderMaskTwoB{} at 2B and \nLadderMaskOneB{} at 1B, so prediction degrades more slowly as more of the scene is masked. Masking is a controlled proxy for road users hidden by others, not a test on real occlusion. In exchange the clips lower Motion cosine at both scales and raise Forward L1 by 0.001 at 2B and Future-frame and Forward L1 by about 0.003 at 1B. On \nHeldoutClips{} clips from \nHeldoutVideos{} videos unused in training, the 1B rows keep every sign and the lead over LoRA holds except on Motion cosine (video-level intervals).
